\section{Research Motivation}
\label{sec:motivation}

This research is primarily motivated by the need to enable autonomous agricultural robots to perform complex manipulation tasks, specifically robotic pruning. Unlike simple navigation, pruning requires a high-fidelity Simultaneous Localisation and Mapping (SLAM) framework that can provide the millimetre-level precision necessary for a robotic arm to interact with thin, overlapping vine structures. Maintaining such accuracy in the presence of dynamic foliage and repetitive, unstructured row patterns remains a critical bottleneck in automating viticulture and horticulture.

The demand for such technology is underscored by the economic landscape of New Zealand's primary industries. As of 2026, horticulture export revenue is forecast to reach a record NZ\$9.2 billion, with wine production constituting NZ\$2.3 billion~\cite{industriesSituationOutlookPrimary2025}. However, the sector faces chronic labour shortages and rising operational costs, with manual pruning alone accounting for approximately 23\% of annual vineyard expenses~\cite{2025MarlboroughVineyard}.

The 2025 Vineyard Monitoring Report published by the Ministry for Primary Industries indicates that during the 2024/2025 season, expenditure on manual pruning increased, while spending on machine pruning declined by 21\%. This shift is attributed to grower concerns regarding vine damage caused by current mechanical stripping machines~\cite{2025MarlboroughVineyard}. These concerns highlight the principal limitation of current commercial machine-based pruning systems: the inability to selectively remove individual canes without compromising cordon or trunk wood. Overcoming these precision challenges would enable robotic automation to provide a scalable response to ongoing workforce constraints, thereby supporting the long-term sustainability of New Zealand’s high-value horticultural crops.

Beyond the immediate application to pruning, this work also contributes to UCVision's broader mission of modelling occluded plant structures and tracking growth across multiple seasons.




